% !PRINTED 499
\weberpage{499}
For if, when $\alpha$ denotes any one of the figures
$0, 1, 2, \ldots n - 1$, we denote by $\alpha_\pi$ and $\alpha_{\pi'}$ the figures into which $\alpha$
passes under $\pi$ and under $\pi'$, then $\alpha_{\pi\pi'}$ is the figure into which $\alpha_\pi$
passes under $\pi'$, or $\alpha$ under $\pi\pi'$. If $\alpha, \beta, \gamma \ldots$ are the figures
from the system $M$, then
\[
\begin{gathered}
\pi'^{-1}(\alpha, \beta) \pi' = (\alpha_{\pi'}, \beta_{\pi'}) \\
\pi\pi'^{-1}(\alpha, \beta) \pi'\pi^{-1} = (\alpha_{\pi'\pi^{-1}}, \beta_{\pi'\pi^{-1}}) = (\gamma, \delta).
\end{gathered}
\]
Now if $\alpha_{\pi'} = \gamma_\pi$ is contained in $M'$, then $\alpha_{\pi'\pi^{-1}} = \gamma$ is contained
in $M$, and consequently $\beta_{\pi'\pi^{-1}} = \delta$ must likewise be contained in $M$,
since the transposition $(\gamma, \delta)$ occurs in $P$. Consequently $\alpha_{\pi'} = \gamma_\pi$
and $\beta_{\pi'} = \delta_\pi$, i.e. if one of the figures from $M''$ is contained in $M'$, then
$M'$ and $M''$ must be identical.
The totality of the figures $0, 1 \ldots n - 1$ therefore falls apart
into systems $M, M', M'' \ldots$ of $m$ figures each, which are permuted imprimitively
by $P$; that is, the group $P$ is imprimitive.
If $n$ is a prime, there are no imprimitive permutation groups;
hence in this case a transitive group which contains a transposition
must be the symmetric group.
In exactly the same way we conclude from Theorem 7:

10. \gesperrt{If a transitive permutation group of
$n$ figures contains a three-term cycle,
then it contains the whole alternating group
or it is imprimitive.}

For if the group $P$ contains the three-term cycles
\[
(0, 1, 2), (0, 1, 3) \ldots (0, 1, m - 1),
\]
but no other with the figures $0, 1$, then by 7.
it contains the whole alternating group of the figures $M = 0, 1, 2 \ldots m - 1$.
Now if $m < n$, no three-term cycle can occur in $P$ which joins one
or two of the figures of $M$ with other figures. For if $P$ contains
e.g. $(0, m, m + 1)$, then, as one sees from the compositions
\[
(0, 2, 1) (0, m, m + 1) (0, 1, 2) = (1, m, m + 1)
\]
it also contains $(1, m, m + 1), \ldots (m - 1, m, m + 1)$ and hence
by 7. the whole alternating group of the figures $0, 1, 2 \ldots m + 1$,
which is contrary to the hypothesis. The $n$ figures of $P$ therefore
fall apart, as before, into several systems $M, M', M'' \ldots$ of imprimitivity. We may add further,
%CONT>>
